% ============================================================
%  Section 1 — Définition du cadre de l'étude
% ============================================================
%
%  Contenu attendu (voir sujet PDF, section "Votre analyse") :
%
%  1.1 Description de la population étudiée
%      - Population cible : étudiant·es P1 EFREI, classes SC5/SC6/SC7 et BDX
%      - Échantillon principal : sondage SC3/SC4 (n = 72 réponses)
%      - Définir : individu, population, échantillon (Déf. 1.1 du cours)
%      - Préciser le mode de recueil (Google Forms, anonyme)
%
%  1.2 Description des variables collectées
%      - Lister chaque variable du sondage ADD_SC3_SC4.csv
%      - Pour chacune : nature (qualitative / quantitative, Déf. 1.3),
%        modalités ou plage de valeurs, rôle dans l'étude
%      - S'appuyer sur data/README.md pour le tableau récapitulatif
%
%  1.3 Discussion des biais potentiels
%      - Biais de sélection (réponses volontaires, classes ciblées)
%      - Biais de déclaration (auto-évaluation, mémoire approximative)
%      - Biais de non-réponse (taux de complétion, questions sensibles)
%      - Discuter la représentativité de l'échantillon
%
% ============================================================

\section{Définition du cadre de l'étude}

\subsection{Population et échantillon}

% TODO : rédiger ici

\subsection{Description des variables collectées}

% TODO : insérer le tableau des variables (voir data/README.md)

\subsection{Biais potentiels}

% TODO : rédiger ici
